\usepackage[utf8]{inputenc}
\usepackage[T1]{fontenc}
% Both languages of the course.  The main language is the edition: the
% driver passes it, \PassOptionsToPackage{main=british}{babel} or
% main=swedish, and didactic's onlyin blocks and \ensv follow it.
% PROJECT: language-specific.
\usepackage[british,english,swedish]{babel}
\usepackage{booktabs}
\usepackage{array}
\usepackage{longtable}
\usepackage{tabularx}

\usepackage[
  natbib,
  style=verbose,
  citestyle=verbose,
  singletitle=false,
  maxbibnames=99,
]{biblatex}
\addbibresource{ltnotes.bib}

\usepackage[all]{foreign}
\renewcommand{\foreignfullfont}{}
\renewcommand{\foreignabbrfont}{}

\usepackage{import}
% beamer loads graphicx itself; memoir does not, and a deck with figures
% needs it in both jobs.
\usepackage{graphicx}

\usepackage[strict]{csquotes}
\usepackage[single]{acro}
% PROJECT: language-specific.  Own-language term first, the English term
% once in parentheses at first use: acro prints "long (short)" on first
% use, so the own-language word has to be the long form and the English
% expansion has to ride along inside it.
% \DeclareAcronym{IDE}{
%   short = IDE,
%   long  = {utvecklingsmiljö
%            (\foreignlanguage{english}{integrated development
%             environment})},
% }

\usepackage{subcaption}

\usepackage{xparse}

\let\email\texttt

% minted v3+ (TeX Live 2024+) auto-detects the output directory; the old
% [outputdir=...] package option is an error there.
\usepackage{minted}
\setminted{autogobble,fontsize=\footnotesize}

% PythonTeX runs the example programs and embeds their real terminal
% sessions (didactic's \runpython[transcript]), so the deck never
% hand-copies its own output: pythontex runs from ltxobj (the outputdir),
% and workingdir=.. moves it up to the deck directory so examples/...
% resolves.  The two editions of a bilingual deck build in different
% output directories (see the Makefile), since PythonTeX's files there
% do not say which job made them.
\usepackage[makestderr]{pythontex}
\setpythontexoutputdir{.}
\setpythontexworkingdir{..}

\usepackage{amsmath}
\usepackage{amssymb}
\usepackage{mathtools}
\usepackage{amsthm}
\usepackage{thmtools,thm-restate}
\usepackage[unq]{unique}

\usepackage[binary-units,parse-numbers=false]{siunitx}

\makeatletter
\@ifclassloaded{beamer}
  {\relax}
  {\usepackage[colorlinks]{hyperref}}
\makeatother
% cleveref takes its language from its own options or the class
% options, not from babel, and its names from the LAST language among
% them: so the driver passes the languages, the edition's last, and
% capitalise for English only.  Load it with none here.  (Without
% capitalise, for Swedish typography, a sentence-initial reference is
% \Cref.)
\usepackage{cleveref}
% Our own reference names, in an \AtBeginDocument hook registered after
% cleveref's own (from the language option) so ours win.  cleveref's
% Swedish names call the appendix type "Appendix", but memoir and babel
% call it "Bilaga"; the prose refers back to numbered lines of a program
% listing as lines ("rader"), not cleveref's "item"/"punkt".  \ensv picks
% the edition's half, so a single-language deck needs no change here.
% capitalise never reaches names set with \crefname, so the English
% halves are capitalised by hand.  Each \ensv is a pair that
% didactic-sync.py stamps; only that script writes the stamps.
\AtBeginDocument{%
  \crefname{appendix}%
    {\ensv{Appendix}{bilaga}} % sync english=2d4b3e swedish=edf9eb
    {\ensv{Appendices}{bilagor}} % sync english=c14d56 swedish=84e46d
  \Crefname{appendix}%
    {\ensv{Appendix}{Bilaga}} % sync english=2d4b3e swedish=2603a2
    {\ensv{Appendices}{Bilagor}} % sync english=c14d56 swedish=9dc348
  \crefname{enumi}%
    {\ensv{Line}{rad}} % sync english=d7852c swedish=9a07c3
    {\ensv{Lines}{raderna}} % sync english=3b26a5 swedish=10766e
  \Crefname{enumi}%
    {\ensv{Line}{Rad}} % sync english=d7852c swedish=aaa6ca
    {\ensv{Lines}{Raderna}} % sync english=3b26a5 swedish=9490da
  \crefname{footnote}%
    {\ensv{Footnote}{fotnot}} % sync english=aa8885 swedish=4c81fc
    {\ensv{Footnotes}{fotnoterna}} % sync english=ad6fb5 swedish=bad1f2
  \Crefname{footnote}%
    {\ensv{Footnote}{Fotnot}} % sync english=aa8885 swedish=a0b3e4
    {\ensv{Footnotes}{Fotnoterna}} % sync english=ad6fb5 swedish=50cd70
}

\usepackage[marginparmargin=outer]{didactic}
% didactic sets footnotes --- and, under the verbose citation style, every
% \autocite --- as margin notes, and LaTeX allocates \marginpar boxes from
% the same pool of eighteen boxes as floats.  A chapter with this many
% margin references empties the pool and dies with "Too many unprocessed
% floats" pages later, in a place that has nothing to do with the cause.
\extrafloats{200}
% memoir places a side caption with \rlap, so it takes part in no
% collision avoidance at all and silently overprints an \ltnote that
% happens to sit at the same height.  Aligning the caption with the BOTTOM
% of the figure moves it below the notes, which anchor where they stand in
% the body text.  (\setsidecappos is memoir's; beamer does not define it.)
\ifdefined\setsidecappos\setsidecappos{b}\fi
\ProvideSemanticEnv{lo}{Learning Objective}
  [style=definition,numbered=yes]
  {LO}{LO}
  {Learning objective}{Learning objectives}
\ProvideTranslation{british}{Learning Objective}{Learning Objective}
\ProvideTranslation{british}{LO}{LO}
\ProvideTranslation{british}{Learning objective}{Learning objective}
\ProvideTranslation{british}{Learning objectives}{Learning objectives}
% PROJECT: language-specific.
\ProvideTranslation{swedish}{Learning Objective}{Lärandemål}
\ProvideTranslation{swedish}{LO}{lm}
\ProvideTranslation{swedish}{Learning objective}{Lärandemål}
\ProvideTranslation{swedish}{Learning objectives}{Lärandemål}

% Paper titles in the generated hit-list tables can carry Greek letters;
% scholar's table export sets them in math mode itself, and this mapping
% stays as a guard for the same characters typed anywhere else.
\DeclareUnicodeCharacter{03C0}{\ensuremath{\pi}}
\DeclareUnicodeCharacter{03BB}{\ensuremath{\lambda}}
\DeclareUnicodeCharacter{03B1}{\ensuremath{\alpha}}
\DeclareUnicodeCharacter{03B2}{\ensuremath{\beta}}
\DeclareUnicodeCharacter{03BC}{\ensuremath{\mu}}
